% titel: Blatt 5, A3: Valenzelektronen in Iodid (Diskussion)

\subsection*{Frage}
Gehören die $4d$-Elektronen von $\mathrm{I^-}$
($[\mathrm{Kr}]\,4d^{10}\,5s^2\,5p^6$) ebenfalls zu den Valenzelektronen?

\subsection*{Antwort: Nein}
Die Valenzelektronen sind definiert als die Elektronen der \textit{äussersten
Schale} (höchste Hauptquantenzahl $n$). Bei Iod ist das die Schale $n=5$
($5s^2\,5p^6$ im Anion). Die $4d^{10}$ liegen in der darunter liegenden Schale
$n=4$:

\begin{itemize}
  \item energetisch deutlich tiefer und fest an den Kern gebunden,
  \item voll besetzt und abgeschlossen,
  \item daher nicht an chemischen Bindungen beteiligt.
\end{itemize}
Sie bilden zusammen mit $[\mathrm{Kr}]$ den inerten Rumpf. Nur die
$5s^2\,5p^6$-Elektronen sind Valenzelektronen; in Verbindungen beteiligt Iod
davon $5s^2\,5p^5$ (bzw.\ bei Oktetterweiterung zusätzlich freie $5d$-Niveaus).

\begin{ergebnis}
Die $4d$-Elektronen in $\mathrm{I^-}$ zählen \textbf{nicht} zu den
Valenzelektronen --- sie gehören zur abgeschlossenen Rumpfschale $n=4$.
\end{ergebnis}

\begin{kontrolle}
Konsistenzprüfung: In $\mathrm{I_2}$ etwa teilt jedes I-Atom nur seine
$5s/5p$-Elektronen in der $\mathrm{I-I}$-Bindung; die $4d^{10}$ bleiben unverändert
Lone Pairs bzw.\ Rumpf --- völlig analog, wie die $3d^{10}$ in $\mathrm{Br^-}$
keine Valenzelektronen sind.
\end{kontrolle}

\begin{merksatz}
Valenzelektronen = Elektronen der \textit{äussersten} Schale. Abgeschlossene
tiefe ($n-1)d$-Unterschalen sind Rumpf --- das unterscheidet Valenz- von
Rumpfelektronen zuverlässig.
\end{merksatz}
